% /Users/pikachu/books/payments-book/src/parts/part5/ch27-ledger.tex
\chapter{Внутренний ledger и событийная модель}
\label{ch:internal-ledger}
\margintoc


\begin{kaobox}[title=Что вы узнаете из этой главы]
  Почему платёжный шлюз не может быть окончательным источником
  правды о деньгах; как auth, capture, refund и dispute превращаются
  во внутренние события и postings; чем pending отличается от available;
  где живут комиссии, резервы и payout schedule; и как связать
  идемпотентность gateway с бухгалтерской консистентностью ledger.
\end{kaobox}

Платёжная система почти всегда выглядит быстрее, чем движутся деньги.
Шлюз отвечает мерчанту за сотни миллисекунд, но в этот момент ещё
не ясно, произойдёт ли capture, когда придёт settlement, какой будет
комиссия эквайера и не превратится ли успешная продажа через месяц
в chargeback. Поэтому между boundary-слоем шлюза и операционной
сверкой нужен ещё один слой: внутренний ledger, который хранит
не просто статусы, а денежные обязательства.

\section{Почему шлюз не может быть системой учёта}
\label{sec:ledger-why}

Шлюз отвечает на вопрос \enquote{что произошло на границе системы}:
пришёл запрос, прошёл risk engine, эквайер вернул approval или decline,
webhook был доставлен мерчанту. Для интеграции этого достаточно,
но для финансового учёта недостаточно по двум причинам.

Во-первых, статус шлюза не равен состоянию денег. Авторизация
\enquote{approved} означает лишь, что эмитент зарезервировал сумму
или обещал её оплатить. До capture и клиринга у мерчанта нет права
на выручку, а у платформы нет основания признавать доход по комиссии.

Во-вторых, один внешний статус может скрывать несколько денежных
последствий. Успешная транзакция в marketplace может породить
отдельные обязательства перед мерчантом, платформой, эквайером
и участниками payout-цепочки. Хранить это в виде одной строки
\codeword{payment.status = succeeded} означает рано или поздно
потерять связь между движением денег и жизненным циклом продукта.

\begin{kaobox}[title={Определение: Internal ledger}, colback=blue!4, colbacktitle=blue!15]
  Internal ledger — внутренняя учётная система платёжного продукта,
  в которой каждая операция представлена как набор денежных postings
  между счетами или балансами. Ledger не заменяет банковский счёт
  и не отменяет внешнюю бухгалтерию, но служит системой истины
  для того, как платёжный продукт понимает собственные обязательства.
\end{kaobox}

Поэтому полезно разделить роли. Шлюз из главы \ref{ch:gateway}
отвечает за приём, маршрутизацию и протокольную надёжность.
Ledger отвечает за состояние денег внутри продукта: что уже
заработано, что только зарезервировано, что удержано под риск,
а что можно выплатить мерчанту.

\section{События: auth, capture, refund, dispute}
\label{sec:ledger-events}

Хороший ledger строится не вокруг CRUD-статусов, а вокруг событий.
Минимальный набор событий для карточного потока выглядит так:
authorization approved, authorization reversed, capture posted,
refund posted, chargeback opened, representment accepted,
chargeback lost, fee accrued.

Каждое такое событие должно быть атомарно записано в event log
и сразу порождать postings в ledger. Тогда продуктовая система
получает две полезные проекции одной и той же реальности:
\textit{что случилось} и \textit{как это изменило деньги}.

Например, авторизация на 5\,000~₽ ещё не создаёт выручку мерчанта,
но создаёт резерв обязательства: система теперь знает, что при
успешном capture эта сумма может перейти из состояния pending
в состояние payable. Capture превращает ожидание в финансовый факт.
Refund делает обратное движение: уменьшает обязательство перед
мерчантом или создаёт отдельную дебиторскую позицию к мерчанту,
если деньги уже были выплачены. Chargeback вообще не является
\enquote{отрицательным refund}: это самостоятельное событие
из внешней схемы со своими сроками, доказательной логикой
и комиссиями (см. главу \ref{ch:disputes}).

\begin{kaobox}[title=На практике]
  Самая частая ошибка в платёжных продуктах — кодировать refund
  и chargeback как один и тот же тип операции со знаком минус.
  Для клиента оба сценария выглядят как возврат денег, но для
  продукта это разные процессы: refund инициирует мерчант, а chargeback
  приходит извне, несёт fee и меняет риск-профиль мерчанта.
\end{kaobox}

\section{Pending, available и postings}
\label{sec:ledger-balances}

Как только события стали первоклассными объектами, следующий шаг —
развести типы балансов. Почти любой платёжный продукт рано или поздно
нуждается минимум в трёх состояниях денег.

\textbf{Pending} — сумма, которая ещё не доступна к выплате,
потому что находится между авторизацией и capture, ждёт клиринга
или удерживается внутри payout window. Pending полезен тем, что
показывает будущую ликвидность, не смешивая её с уже доступными
средствами.

\textbf{Available} — сумма, которой платформа уже готова
распоряжаться: перечислить мерчанту, зачесть в счёт комиссии
или использовать для внутреннего неттинга.

\textbf{Reserve / holdback} — сумма, изъятая из available
по риск-правилу: rolling reserve для high-risk мерчанта,
удержание под ожидаемые chargebacks, резерв на refund period.

С инженерной точки зрения это означает, что одна внешняя транзакция
обычно порождает несколько postings. Capture может дебетовать
счёт \enquote{acquirer receivable} и кредитовать \enquote{merchant pending}.
После получения settlement ledger переводит сумму из
\enquote{merchant pending} в \enquote{merchant available},
одновременно начисляя processing fee на отдельный счёт платформы.

Если эти переходы не моделировать явно, продукт быстро начинает
путать \enquote{деньги есть} с \enquote{деньги когда-нибудь придут}.
Ровно из этого обычно рождаются загадочные расхождения между
дашбордом мерчанта, payout-отчётом и банковской выпиской.

\section{Комиссии, резервы и payout schedule}
\label{sec:ledger-fees}

Ledger становится по-настоящему полезным в тот момент, когда
в нём появляются не только основные суммы платежей, но и все
производные денежные эффекты.

Комиссия платформы редко возникает в один момент с authorizaton.
Обычно она признаётся при capture или settlement, когда уже ясно,
что операция действительно состоялась. Эквайерские сборы, сетевые fee,
FX markup и комиссии за chargeback могут приходить с иным таймингом,
чем основная сумма транзакции. Значит, ledger должен поддерживать
accrual-логику: сначала признать ожидаемый fee, потом скорректировать
его по фактическому settlement-файлу.

То же относится к резервам. Rolling reserve на 10\% означает не
просто \enquote{не заплати мерчанту всё сразу}, а конкретное правило
перевода части available в отдельный баланс с собственным сроком
разморозки. Если это правило не выражено в ledger, его начинают
повторять вручную в payout-job, backoffice и support-инструментах,
после чего система теряет единый источник истины.

Payout schedule тоже лучше моделировать в ledger, а не в коде
экспортёра банковских файлов. Выплата раз в день, T+2, weekly payout
по пятницам, ручная задержка high-risk мерчанта — всё это правила
перевода между внутренними балансами, а не просто cron-выражения.

\section{Идемпотентность между gateway и ledger}
\label{sec:ledger-idempotency}

Шлюз и ledger живут в разных временных режимах. Шлюз должен быстро
вернуть API-ответ мерчанту. Ledger должен гарантировать, что событие
будет отражено ровно один раз, даже если внешний мир несколько раз
повторил один и тот же сигнал.

Практически это означает две границы идемпотентности. Первая —
на входе в gateway: мерчант повторяет запрос с тем же
\codeword{idempotency-key}, а шлюз не должен создавать вторую
авторизацию. Вторая — между gateway и ledger: webhook от эквайера,
повторный callback после timeout или переигранный batch не должны
создать дубль postings.

Хороший паттерн здесь такой:
\begin{itemize}
  \item gateway присваивает операции собственный стабильный
        идентификатор;
  \item все внешние сообщения связываются с ним через mapping
        STAN / RRN / processor reference;
  \item ledger принимает только нормализованные внутренние события
        с уникальным \codeword{event_id};
  \item повторная доставка события возвращает прежний результат,
        не порождая новых postings.
\end{itemize}
Именно здесь удобно развести \enquote{event store} и \enquote{ledger postings}.
Сначала система фиксирует факт получения события от внешнего мира,
затем атомарно применяет его к денежным балансам. Если между этими
шагами произойдёт сбой, операцию можно безопасно доиграть.

\section{Передача в reconciliation}
\label{sec:ledger-handoff}

Ledger не отменяет reconciliation, а делает его осмысленным.
Сверка из главы \ref{ch:reconciliation} отвечает на вопрос,
совпадает ли внутренняя картина продукта с внешними файлами
и банковскими движениями. Чтобы этот вопрос вообще можно было
задать, продукт сначала должен знать собственную внутреннюю картину.

Удобно мыслить handoff так:
\begin{itemize}
  \item gateway фиксирует протокольный факт авторизации;
  \item ledger превращает его в денежное состояние;
  \item reconciliation сверяет это состояние с TC33, IPM,
        settlement-отчётами и bank statement.
\end{itemize}
Если ledger отсутствует, reconciliation пытается напрямую
сопоставить внешние файлы со статусами шлюза, и команда получает
типичный хаос: \enquote{вроде approved, но непонятно, это уже доход,
ожидаемая выплата или просто авторизационный hold}. Ledger снимает
эту неоднозначность ещё до начала сверки.

\begin{kaobox}[title=На практике]
  Самый полезный артефакт для операционной команды — таблица
  соответствия между внешними объектами и внутренними счетами:
  authorization ID, processor reference, merchant order ID,
  payout ID, dispute case ID. Без неё даже хороший ledger превращается
  в чёрный ящик, который невозможно сверить с реальным миром.
\end{kaobox}

\begin{chaptersummary}
\begin{itemize}
  \item Шлюз отвечает за boundary-слой и протокол, но не может
        быть окончательной системой учёта денег.
  \item Базовые платёжные события — auth, capture, refund, dispute,
        fees — должны сначала фиксироваться как события, а затем
        превращаться в postings.
  \item Разделение pending, available и reserve защищает продукт
        от смешения будущих, доступных и удержанных средств.
  \item Комиссии, rolling reserve и payout schedule лучше моделировать
        как правила ledger, а не как разрозненную логику в backoffice.
  \item Идемпотентность нужна не только на API-входе, но и на границе
        между gateway и ledger: повторная доставка события не должна
        создавать дубль postings.
  \item Reconciliation начинается после ledger: только имея внутреннюю
        денежную картину, можно осмысленно сверять её с клирингом,
        settlement и банковской выпиской.
\end{itemize}
\end{chaptersummary}

